% ============================================================
\section{Experiments}
\label{sec:experiments}
% ============================================================

% ============================================================
\subsection{Experimental Setup}
\label{sec:experiments:setup}
% ============================================================

\textbf{Dataset.}
We use \dataset{}~\cite{neto2023ciciot}, a publicly available benchmark
of labelled IoT network flows captured from 105 IoT devices spanning
33 attack categories.  We extract 12 network-flow features and create
128-timestep sliding windows with stride 128, using a stratified 80/20
train/test split.

\textbf{Synthetic stealth-controlled negatives.}
We generate 4 attack types (Eqs.~\ref{eq:slow}--\ref{eq:proto})
at 3 stealth levels (85\%, 90\%, 95\%), yielding 12 evaluation
conditions.  Protocol constraints (Modbus, MQTT, CoAP) are enforced
at the \emph{moderate} strictness tier.
We compare two negative-generation pipelines: (1)~analytical
closed-form perturbations applied to real benign samples (condition~D),
and (2)~the same perturbations applied to \emph{learned} base samples
from Diffusion-TS~\cite{zhang2024diffusionts} with post-hoc statistical
guidance (condition D-DiffTS; Section~\ref{sec:method:diffts}).
This controlled comparison isolates the effect of the base distribution
(real vs.\ generated) while holding the attack perturbation constant.

\textbf{Threshold calibration.}
We set the anomaly threshold at the 95th percentile of a benign-only
held-out set (20\% of benign samples), targeting FPR\,$\leq$\,0.05.
This protocol is applied uniformly to all methods; only benign data
inform the threshold (details in Appendix~\ref{app:threshold}).

\textbf{Baselines.}
We compare against nine methods (5 traditional, 4 deep-learning; see
Appendix~\ref{app:baseline_tuning}).
All baselines are \textbf{unsupervised} (benign-only training);
\ours{} conditions B--D are \textbf{supervised} (labeled negatives
during training; see Appendix~\ref{app:supervision}).

\textbf{Metrics.}
We report F1, FPR, and ROC-AUC.  All experiments use 10 random seeds
except leave-one-out (3 seeds, due to combinatorial cost).

% ============================================================
\subsection{Main Results}
\label{sec:experiments:main}
% ============================================================

Table~\ref{tab:main} compares all methods across three synthetic stealth levels.

% Table 1: Main Results — baselines on combined test set + DiffIDS per stealth level
% Results from run_all_baselines.py (baselines) and run_ablation.py --condition d (DiffIDS).

\begin{table}[t]
\centering
\caption{
  Detection performance.  Baselines evaluated on the combined test set (all stealth
  levels mixed).  \ours{} reported at each stealth level separately, demonstrating
  robustness against increasingly stealthy attacks.
  \textbf{Bold}: best per column.  \underline{Underline}: second best.
  $\dagger$ marks our proposed system.
}
\label{tab:main}
\resizebox{\textwidth}{!}{%
\begin{tabular}{lcccccccc}
\toprule
& \multicolumn{2}{c}{Combined Test Set} & \multicolumn{2}{c}{Stealth 85\%} & \multicolumn{2}{c}{Stealth 90\%} & \multicolumn{2}{c}{Stealth 95\%} \\
\cmidrule(lr){2-3} \cmidrule(lr){4-5} \cmidrule(lr){6-7} \cmidrule(lr){8-9}
Method & F1 & FPR & F1 & FPR & F1 & FPR & F1 & FPR \\
\midrule
\multicolumn{9}{l}{\textit{Traditional baselines}} \\
Threshold IDS       & 0.878 & 0.600 & \multicolumn{6}{c}{\textit{(not stealth-level specific)}} \\
Signature IDS       & 0.000 & 0.000 & & & & & & \\
Statistical IDS     & 0.889 & 1.000 & & & & & & \\
Isolation Forest    & 0.302 & 0.500 & & & & & & \\
Ensemble            & \underline{0.857} & \underline{0.400} & & & & & & \\
\midrule
\multicolumn{9}{l}{\textit{Deep-learning baselines}} \\
USAD                & 0.613 & 0.300 & & & & & & \\
TranAD              & 0.423 & 0.100 & & & & & & \\
Anomaly Transformer & 0.696 & 0.500 & & & & & & \\
PatchTST-Anomaly    & 0.567 & 0.300 & & & & & & \\
\midrule
\multicolumn{9}{l}{\textit{Ours (Moirai-based)}} \\
\ours{} zero-shot (A) & 0.860 & 0.933 & 0.691 & 0.767 & 0.682 & 0.933 & 0.706 & 0.833 \\
\ours{}$^\dagger$ (D) & \textbf{0.877} & \textbf{0.500} & \textbf{0.818} & \textbf{0.300} & \textbf{0.763} & \textbf{0.567} & \textbf{0.732} & \textbf{0.500} \\
\bottomrule
\end{tabular}
}
\end{table}


At stealth-95 with benign-calibrated thresholds, Signature IDS and
Statistical IDS are degenerate (FPR$\geq$0.87; Table~\ref{tab:main},
$\ddagger$ rows).  Among calibrated methods, the Ensemble achieves
F1$=$0.136 (FPR$=$0.042).
\ours{} condition~C (F1$=$0.262, FPR$=$0.068) outperforms all calibrated
baselines at stealth-95.
While this is the highest calibrated performance among all methods,
an absolute F1 of 0.262 means roughly one in four stealth-95 attacks
is detected under operational thresholds---a diagnostic rather than
operational result (see Section~\ref{sec:analysis} for deployment guidance).
Zero-shot A (F1$=$0.137, FPR$=$0.063) performs at chance level, consistent
with its near-random AUC (0.481$\pm$0.042).

% ============================================================
\subsection{Ablation Study}
\label{sec:experiments:ablation}
% ============================================================

Table~\ref{tab:ablation} ablates the five components of \ours{} (200
samples, 5 epochs, 10 seeds, benign-calibrated threshold).

% Table 2: Ablation Study
% Auto-generated by scripts/generate_paper_figures.py after experiments complete.

\begin{table}[h]
\centering
\caption{
  Ablation study on stealth-95 (hardest) and all stealth levels combined.
  Each condition adds one component to the previous.
  $\Delta$F1 is relative to condition A (zero-shot baseline).
}
\label{tab:ablation}
\begin{tabular}{llcccc}
\toprule
& & \multicolumn{2}{c}{Stealth 95\%} & \multicolumn{2}{c}{All Stealth} \\
\cmidrule(lr){3-4} \cmidrule(lr){5-6}
Cond. & Description & F1 & FPR & F1 & FPR \\
\midrule
A & Moirai zero-shot          & -- & -- & -- & -- \\
B & + Fine-tune (NLL only)    & -- & -- & -- & -- \\
C & + SupCon (real attacks)   & -- & -- & -- & -- \\
D & + Hard negatives (\ours{}) & \textbf{--} & \textbf{--} & \textbf{--} & \textbf{--} \\
E & D w/o protocol constraints & -- & -- & -- & -- \\
\bottomrule
\end{tabular}
\end{table}


\textbf{Fine-tuning improves over zero-shot, but with high variance.}
We report bootstrap 95\% CIs, Cohen's $d$ effect sizes, and $p$-values for
key comparisons in Appendix~\ref{app:statistics}.
The best Moirai condition (C, AUC$=$0.629) significantly improves $+$0.148
over zero-shot A (AUC$=$0.481; $p{<}0.001$, $d{=}4.45$);
B (0.556) also significantly exceeds A ($p{<}0.001$, $d{=}2.29$).
The from-scratch CNN (AUC$=$0.580) is significantly below C ($p{<}0.001$);
a CNN variant with Gaussian NLL (CNN-NLL, AUC$=$0.598) confirms the
loss function is not the explanation ($p{=}0.24$ vs.\ CNN-MSE).

\textbf{Seed variance.}
At 10 seeds, Moirai variance stabilises substantially (C: $\pm$0.022 vs.\
$\pm$0.122 at 5 seeds); the $5.5\times$ reduction exceeds the expected
$\sqrt{2}$ factor (post-hoc observation: seeds 789 and 1234 appear to have
depressed the 5-seed mean; this was identified after unblinding and should
be treated as exploratory).  CNN variance is comparable ($\pm$0.023).
All main results and scaling experiments use 10 seeds; leave-one-out uses
3 seeds due to combinatorial cost (see Appendix~\ref{app:cnn_detail}).
Condition~C$'$ (balanced stealth-controlled negatives, 1:1) disentangles
negative type from imbalance: Gaussian-noise negatives outperform
stealth-controlled negatives by $+$0.019 AUC
(C vs.\ C$'$; $p{=}0.070$, $d{=}0.82$), and the 1:12 imbalance adds a
significant $-$0.054 penalty (C$'$ vs.\ D; $p{<}0.001$, $d{=}2.33$).
See Appendix~\ref{app:cprime} for the full decomposition.
Conditions E, E$'$, and E$''$ confirm that the constraint validator and
retry mechanism have negligible effect on the analytical pathway
(Appendix~\ref{app:eee}; embedding visualization in Appendix~\ref{app:embeddings}).
SupCon hyperparameters ($\lambda$, $\tau$) have no detectable effect at
this scale (Appendix~\ref{app:hparams}).

% Protocol constraint compliance (100% at all tiers) is reported in
% Appendix~\ref{app:constraints}; this is expected for the analytical pathway.
\label{sec:experiments:constraints}

% ============================================================
\subsection{Generalization to Unseen Attack Patterns}
\label{sec:experiments:loo}
% ============================================================

A leave-one-out experiment trains condition~D on three attack types and
evaluates on the held-out fourth (3 seeds, 30 eval samples per stealth level).
All four folds show positive transfer: fine-tuning improves over zero-shot
for every held-out attack type (Table~\ref{tab:loo}, Appendix~\ref{app:loo_detail}).
The largest gain is for \texttt{lotl\_mimicry}, suggesting
that low-amplitude mimicry benefits most from boundary-shaping by
the other archetypes.
\textbf{Important caveats:} this is a limited generalisation test.
With $k{=}4$ hand-designed archetypes, each fold trains on only
three types from a human-selected taxonomy.
With only 3 seeds and 30 eval samples per stealth level, per-fold
standard deviations ($\pm$0.05 to $\pm$0.17) approach the signal
magnitude; the wide confidence intervals mean these results are
\emph{directional signals of positive transfer}, not a comprehensive
generalisation guarantee.
Testing against held-out archetype families outside the original taxonomy
would strengthen this claim.

% ============================================================
\subsection{Scaling Experiment}
\label{sec:experiments:scaling}
% ============================================================

Table~\ref{tab:scaling} compares stealth-95 metrics at three scales
(200/5, 500/10, 1000/20) under NLL-only early stopping (10 seeds),
plus encoder-freezing and L2-SP regularisation diagnostics at 1000/20.

% Table: Scaling Experiment
% Seven regimes: 200/5 (original), 500/10 NLL-ES, 1000/20 NLL-only ES,
%                2000/30 NLL-only ES, 1000/20 frozen encoder, 1000/20 partial freeze
% All stealth-95, benign-calibrated threshold, 5 seeds

\begin{table}[h]
\centering
\caption{
  Scale degradation and encoder-freezing experiment (stealth-95, 5 seeds,
  benign-calibrated threshold, NLL-only early stopping).
  \textbf{Key findings}:
  (1)~D-DiffTS improves from 200/5 to 500/10 (0.621$\to$0.680), then stabilises;
  C/D degrade.
  (2)~Frozen D-DiffTS achieves best AUC (0.728, $+$0.146 over frozen C).
  2000/30 column shows only C, D, D-DiffTS; B, C$'$, CNN not run at this scale.
  \emph{Note}: 200/5 uses 5 seeds for consistency; authoritative 10-seed values in Table~\ref{tab:ablation}.
}
\label{tab:scaling}
\resizebox{\textwidth}{!}{%
\begin{tabular}{@{}llcccccccccc@{}}
\toprule
& & \multicolumn{2}{c}{200\,/\,5} & \multicolumn{2}{c}{500\,/\,10} & \multicolumn{2}{c}{1000\,/\,20} & \multicolumn{2}{c}{2000\,/\,30} & \multicolumn{2}{c}{Frozen 1000/20} \\
\cmidrule(lr){3-4} \cmidrule(lr){5-6} \cmidrule(lr){7-8} \cmidrule(lr){9-10} \cmidrule(lr){11-12}
Cond. & Description & AUC & F1 & AUC & F1 & AUC & F1 & AUC & F1 & AUC & F1 \\
\midrule
B      & NLL only                & $.481_{\pm.127}$ & $.131_{\pm.061}$ & $.538_{\pm.038}$ & $.250_{\pm.073}$ & $.521_{\pm.024}$ & $.177_{\pm.031}$ & --- & --- & $.606_{\pm.114}$ & $.130_{\pm.094}$ \\
C      & SupCon $+$ Gaussian     & $.563_{\pm.122}$ & $.250_{\pm.135}$ & $.592_{\pm.063}$ & $.280_{\pm.101}$ & $.508_{\pm.041}$ & $.175_{\pm.048}$ & $.506_{\pm.140}$ & $.192_{\pm.117}$ & $.582_{\pm.099}$ & $.148_{\pm.097}$ \\
C$'$   & SupCon $+$ hard neg 1:1 & $.534_{\pm.125}$ & $.183_{\pm.084}$ & $.580_{\pm.064}$ & $.265_{\pm.121}$ & $.508_{\pm.025}$ & $.158_{\pm.040}$ & --- & --- & \multicolumn{2}{c}{---} \\
CNN    & 1D-CNN from scratch     & $.597_{\pm.096}$ & $.236_{\pm.069}$ & \multicolumn{2}{c}{---} & $.513_{\pm.029}$ & $.207_{\pm.017}$ & --- & --- & \multicolumn{2}{c}{(n/a)} \\
D      & Hard neg 1:12 (\ours{}) & $.479_{\pm.123}$ & $.137_{\pm.086}$ & $.547_{\pm.044}$ & $.250_{\pm.073}$ & $.506_{\pm.029}$ & $.164_{\pm.059}$ & $.488_{\pm.070}$ & $.164_{\pm.117}$ & $.573_{\pm.097}$ & $.144_{\pm.081}$ \\
\addlinespace[2pt]
D-DiffTS & Learned neg (Diffusion-TS) & $.621_{\pm.056}$ & $.213_{\pm.118}$ & $\mathbf{.680}_{\pm.064}$ & $.297_{\pm.104}$ & $.576_{\pm.045}$ & $.253_{\pm.109}$ & $.578_{\pm.117}$ & $.257_{\pm.102}$ & $\mathbf{.728}_{\pm.080}$ & $.249_{\pm.125}$ \\
\bottomrule
\multicolumn{12}{@{}l}{\scriptsize ``---'': not run. ``(n/a)'': CNN has no pre-training. D-DiffTS 200/5: 10 seeds; others: 5 seeds.} \\
\end{tabular}%
}
\end{table}


\textbf{Monotonic degradation with scale (10 seeds).}
With 10-seed evaluation, all conditions degrade monotonically from
200/5 onward: C drops from 0.629 to 0.581 (500/10) to 0.581 (1000/20);
D from 0.556 to 0.536 to 0.557; B from 0.556 to 0.535 to 0.550.
The earlier appearance of 500/10 improvement (5-seed data) was a
small-sample artefact: the 5-seed 200/5 means were underestimates
(B: 0.481 vs.\ 10-seed 0.556; C: 0.563 vs.\ 0.629).
The C$>$D gap narrows monotonically from $+$0.073 (200/5, $p{<}0.001$)
to $+$0.045 (500/10, $p{=}0.027$)
to $+$0.024 (1000/20, n.s., $p{=}0.216$)---the advantage of easy
off-manifold negatives diminishes as catastrophic forgetting erases
the signal that SupCon operates on.
Variance increases with scale (C: $\pm$0.022 at 200/5, $\pm$0.075 at
500/10, $\pm$0.082 at 1000/20), reflecting seed-dependent forgetting
trajectories.

\textbf{Learned negatives resist forgetting.}
D-DiffTS (5 seeds) improves from 0.621 (200/5) to 0.680 (500/10), the best
unfrozen AUC at any scale.  Unlike Gaussian negatives, which degrade from
0.629 (200/5) to 0.581 (500/10, 10 seeds), D-DiffTS at 1000/20
(0.576) exceeds all other unfrozen conditions.
At 2000/30, D-DiffTS stabilises (0.578$\pm$0.117), while C and D
plateau/decline (C$=$0.506, D$=$0.488; 5-seed 2000/30 values).
Over the full training range (200/5 to 2000/30), D-DiffTS loses only
$-$6.9\% AUC (0.621$\to$0.578), compared to $-$19.6\% for
Gaussian (0.629$\to$0.506).
NLL-only early stopping is essential: under total-loss ES, decreasing SupCon loss masks increasing NLL, causing late checkpoint selection that compounds forgetting (Appendix~\ref{app:nll_es}).

\textbf{Encoder freezing confirms catastrophic forgetting.}
Freezing the encoder at 1000/20 (10 seeds) yields
B$=$C$=$D$=$0.549$\pm$0.109---identical across conditions because
evaluation uses NLL scores from the frozen encoder, and only the
(unused-at-inference) projection head differs.
This result still confirms catastrophic forgetting: frozen 0.549
vs.\ the best unfrozen condition C at 200/5 (0.629) shows that
pre-trained features, when preserved, provide a detection floor
that extended fine-tuning erodes.
However, frozen AUC does \emph{not} exceed the unfrozen 1000/20
conditions (C$=$0.581, B$=$0.550, D$=$0.557), indicating that
the projection head alone cannot recover the discrimination lost
by not adapting the encoder.
Frozen-encoder F1 remains low (0.152$\pm$0.077) under calibrated
thresholds, as the NLL-based scorer cannot leverage the contrastive
signal (details in Appendix~\ref{app:freezing_details}).
Partial freezing yields modest further gains
(Appendix~\ref{app:partial_freeze}).

\textbf{Controlled evaluation reveals D-DiffTS advantage is an evaluation-distribution artifact.}
Frozen D-DiffTS previously reported AUC$=$0.728 (5 seeds) using
DiffTS-generated evaluation data; unfrozen D-DiffTS at 1000/20 achieves
0.576 on the same DiffTS evaluation set.
To decouple training data from evaluation distribution, we re-evaluate
\emph{both} frozen and unfrozen D-DiffTS on the same analytical negatives
used for B/C/D (5 seeds, 1000/20; Table~\ref{tab:scaling};
Appendix~\ref{app:controlled_frozen}).
Frozen D-DiffTS controlled: AUC$=$0.464$\pm$0.060.
Unfrozen D-DiffTS controlled: AUC$=$0.463$\pm$0.073.
Both fall below zero-shot A (0.481) and below all unfrozen B/C/D conditions.
This confirms that the D-DiffTS advantage---both frozen and unfrozen---reflects
the frozen or adapted encoder trivially separating DiffTS negatives
(different feature-space region) rather than any genuine improvement from
learned-negative training when measured on a common evaluation set.
Full methodology and results appear in Appendix~\ref{app:controlled_frozen}.

\textbf{Learning-rate ablation rules out LR misconfiguration.}
To test whether the 1000/20 degradation is merely due to an excessive learning
rate, we re-run B, C, and D at $\text{lr}{\in}\{10^{-5}, 10^{-6}\}$
(Table~\ref{tab:lr_ablation}).
At $10^{-5}$, condition~C partially recovers
(0.557$\pm$0.103 vs.\ $10^{-4}$: 0.508$\pm$0.041) but remains above frozen
(0.549$\pm$0.109); conditions B and D do not benefit.
At $10^{-6}$, all conditions underfit (B: 0.465, C: 0.471, D: 0.465),
confirming a narrow LR window rather than a simple ``lower is better'' fix.
LoRA fine-tuning (rank$=$8, $\alpha{=}$16) on \texttt{q\_proj}, \texttt{v\_proj},
\texttt{out\_proj} also underperforms full freezing (C: 0.484$\pm$0.108,
D: 0.484$\pm$0.133), suggesting the low-rank adapters lack sufficient capacity
for this task.
Together, these results provide converging evidence that encoder-weight
corruption---not hyperparameter misconfiguration---drives the scaling
degradation (full results in Table~\ref{tab:lr_ablation},
Appendix~\ref{app:lr_ablation}).

\textbf{Analytical vs.\ learned negatives.}
To test whether the analytical perturbations (Eqs.~\ref{eq:slow}--\ref{eq:proto})
limit stealth-controlled negative quality, we train a
Diffusion-TS~\cite{zhang2024diffusionts} generative model on CICIoT2023 benign
traffic (Section~\ref{sec:method:diffts}) and generate \emph{learned}
negatives with post-hoc statistical guidance.
D-DiffTS is included as a full row in Tables~\ref{tab:ablation}
and~\ref{tab:scaling} at all three training scales, with frozen-encoder
results at 1000/20.
At 200/5 (10 seeds), D-DiffTS (AUC$=$0.621$\pm$0.056) matches Gaussian C
(0.629$\pm$0.022) on their respective evaluation sets.
At 500/10, D-DiffTS achieves 0.680 and at 1000/20 (frozen), 0.728---but
\emph{both} figures use DiffTS-generated evaluation data, not the analytical
negatives used by B/C/D.
As shown above, controlled cross-evaluation collapses both to $\approx$0.464,
confirming these figures do not reflect a genuine learned-negative advantage.

\textbf{Is the foundation model necessary for learned negatives?}
To test whether DiffTS benefits require Moirai's pre-trained features, we
train from-scratch CNN and CNN-NLL with DiffTS negatives (5 seeds, 200/5).
CNN$+$DiffTS achieves AUC$=$0.514$\pm$0.095 and CNN-NLL$+$DiffTS achieves
0.512$\pm$0.108 at stealth-95, both \emph{below} the same CNNs with
analytical negatives (CNN: 0.580, CNN-NLL: 0.598; Table~\ref{tab:ablation}).
DiffTS negatives \emph{hurt} the from-scratch CNN, confirming that the
D-DiffTS advantage is specific to the foundation model: Moirai's pre-trained
features can leverage the distributional richness of learned negatives,
while a from-scratch CNN lacks the representational capacity to benefit.

\textbf{Generator ablation.}
To isolate the source of any D-DiffTS improvement, we compare three
conditions at 200/5 (5 seeds, Table~\ref{tab:gen_ablation}):
(1)~D-DiffTS (full: learned base $+$ analytical perturbation $+$
post-hoc guidance);
(2)~D-DiffTS-noguide (learned base $+$ perturbation, no guidance);
(3)~D-analytical (existing condition~D: analytical base $+$ perturbation).

% Table: Generator Ablation (D-DiffTS components)
% All conditions: 200/5, 5 seeds, stealth-95, benign-calibrated threshold

\begin{table}[h]
\centering
\caption{
  Generator ablation at 200/5 (stealth-95, 5 seeds).
  Isolates the contribution of the learned base distribution vs.\
  post-hoc statistical guidance.
  D-analytical uses closed-form benign samples;
  D-DiffTS variants use Diffusion-TS-generated benign samples.
}
\label{tab:gen_ablation}
\begin{tabular}{@{}llcc@{}}
\toprule
Condition & Base $+$ Guidance & AUC & F1 \\
\midrule
D-analytical      & Analytical $+$ none     & $.556_{\pm.021}$ & $.176_{\pm.032}$ \\
D-DiffTS-noguide  & Learned $+$ none        & $.465_{\pm.086}$ & $.137_{\pm.057}$ \\
D-DiffTS (full)   & Learned $+$ guidance    & $.621_{\pm.056}$ & $.213_{\pm.118}$ \\
\bottomrule
\end{tabular}
\end{table}


% ============================================================
\subsection{External Validation on N-BaIoT}
\label{sec:experiments:nbaiot}
% ============================================================

To address circular evaluation (Limitation~5), we evaluate on
N-BaIoT~\cite{meidan2018n}, an independent IoT botnet benchmark.
Zero-shot transfer from CICIoT2023 fails (AUC$<$0.30), as expected
given the distribution shift (Appendix~\ref{app:extended_nbaiot}).

\textbf{Device~4 (1{,}096 windows): strongest cross-dataset result.}
On device~4 (Philips B120N10, 1{,}096 training windows, 5 epochs),
D-DiffTS learned negatives achieve \textbf{AUC$=$0.939$\pm$0.014} (3 seeds),
slightly exceeding Gaussian on the same device (0.929$\pm$0.012).
This is the strongest cross-dataset result, confirming that learned
negatives match or exceed Gaussian when adequate training data ($\geq$500
windows) is available.

\textbf{Device~1 (310 windows): data-size limitation.}
On device~1 (Danmini doorbell, 310 training windows), Gaussian-noise
negatives achieve AUC$=$\nbaiotAUC{} (3 seeds), with 8/10 attack
variants above 0.83.
D-DiffTS on this device achieves only AUC$=$0.660$\pm$0.427---the training
set is insufficient for DiffTS to learn a robust benign distribution.
The device~4 result directly confirms the data-size hypothesis.

\textbf{Caveat: N-BaIoT success uses Gaussian negatives, not stealth-controlled.}
The N-BaIoT fine-tuning uses Gaussian-noise negatives, not the
stealth-controlled negatives from the CICIoT2023 pipeline.
This validates the \emph{training recipe} (NLL$+$SupCon) and the
\emph{framework} (Moirai fine-tuning), but does not directly validate
the stealth-controlled negative generation methodology on an
external dataset.

Against unsupervised baselines (Table~\ref{tab:nbaiot_baselines},
Appendix~\ref{app:extended_nbaiot}), Isolation Forest (AUC$=$0.991) leads
with 72 hand-engineered features; \ours{} (AUC$=$\nbaiotAUC{}) matches
One-Class SVM (0.931) using a general-purpose backbone.
